\subsection{Multi-minor locality tournament, W33 routing, distance-two decoding, and hardware capture request}

The $[[20,7,2]]_3\to[[240,81,3]]_3$ symplectic adapter has been extended from a
single fixed-minor support optimization to a finite, explicitly scored tournament
of distinct rank-preserving minors.  For every displayed candidate minor, its
twenty columns are frozen so that $\operatorname{rank}A=20$ is structural.  Every
remaining column is then solved by deterministic CP-SAT in two lexicographic
stages,
\[
  \min |\operatorname{supp}(a_j)|,
  \qquad
  \min\sum_{r\in\operatorname{supp}(a_j)}d_{L(W33)}(e_r,e_{\mathrm{anchor}(r)}),
\]
subject to $Ha_j=t_j$ over $\mathbf F_3$.  For the resulting $A$, each row of
$B$ is optimized by the same exact support-first/locality-second rule subject to
$Ab_i^T=e_i$.  Candidate classes are compared by total $A+B$ support, anchor
locality, maximum row radius, and W33 line-graph spanning-tree cost.  The winner
is therefore exact inside each displayed fixed-minor class and best among the
finite candidate family; no claim is made that all $\binom{240}{20}$ possible
rank minors have been exhausted.

The selected pair $(A,B)$ is lowered all the way to an explicit topological
nearest-neighbour circuit.  The $240$ physical carrier qutrits are the $240$
edges of $W(3,3)$, and a primitive two-qutrit interaction is declared local only
when the two carrier edges share a W33 point.  Remote ternary
\texttt{SUM\_ALPHA} and \texttt{SWAP} macros are replaced by shortest-path swap
networks in the $240$-vertex edge line graph.  Each replacement is replayed over
$\mathbf F_3$ and required to equal the original linear Clifford action.  Route
supports that overlap are serialized, while disjoint supports remain parallel.
The certificate reports topological depth, shortest-path lengths, physical-edge
load, point-coupler load, and a content hash of the resulting local schedule.
This is a topological locality theorem for the explicit W33 interaction graph,
not yet a calibrated optical loss/crosstalk or fault-spread theorem.

The mapped CSS syndrome layer is now explicit as well.  The external code uses
rank-two $X$ checks $H_0$ and rank-eleven $Z$ checks $H^\perp$, hence
$20-2-11=7$ logical qutrits.  Under $(x,z)\mapsto(xA,zB)$ the mapped checks are
$H_0A$ and $H^\perp B$.  Thirteen weighted-check ancillas are scheduled into
72-tick packet microframes with no data-edge or check-ancilla collision inside a
round.  The distance-two boundary is kept exact: all $20\times8=160$
nonidentity single-qutrit Pauli errors are required to produce nonzero syndrome,
but an unknown-location syndrome is corrected only when unique.  Ambiguous
syndromes are detected and refused.  If the damaged coordinate is known as an
erasure, all nine Pauli possibilities at that coordinate are distinguished and
the corresponding correction is emitted as the W33 Pauli frame
\[
  X:\;-aA_q,
  \qquad
  Z:\;-bB_q.
\]
Thus the decoder closes exact detection and one-known-erasure recovery without
misrepresenting a $d=2$ code as an arbitrary single-error-correcting code.

A parameterized noise experiment is tied to the routed circuit while remaining
separate from physical fault-tolerance admission.  If $E_i$ is the conservative
routed-program exposure count of input qutrit $i$, the stated phenomenological
model uses
\[
  p_i=1-(1-p_{\rm gate})^{E_i}
\]
and distributes each nonidentity qutrit Pauli uniformly over its eight choices.
All ${20\choose2}8^2=12160$ weight-two Pauli patterns are exhaustively classified
as detected, stabilizer-trivial, or zero-syndrome logical.  Weight-zero,
weight-one, and weight-two probability are evaluated exactly at each grid point;
all weight-$\ge3$ mass is retained adversarially rather than discarded.  The
reported interval is a conditional \emph{block} logical-failure interval.  No
unsupported division by the seven encoded logical coordinates is used, and a
grid point is called below bare $p$ only if even the adversarial block upper
bound is below $p_{\rm gate}$.  This is a conservative routed-exposure
pseudothreshold experiment, not a calibrated photonic threshold.

The admission policy now distinguishes these evidence classes explicitly:
\[
\begin{array}{rcl}
\text{exact symplectic embedding} &:& \text{implemented},\\
\text{W33 topological route compiler} &:& \text{implemented},\\
\text{mapped packet decoder} &:& \text{implemented},\\
\text{routed-exposure block pseudothreshold} &:& \text{implemented},\\
\text{calibrated optical locality/noise propagation} &:& \text{open},\\
\text{physical mapped threshold certificate} &:& \text{open},\\
\text{canonical }[[66,8,3]]_3\text{ substrate witness} &:& \text{open}.
\end{array}
\]
Consequently fault-tolerant magic-resource admission remains fail closed.

Finally, the Holotrade hardware lane was changed from a manual-only workflow to
an explicit execution-request mechanism.  A request committed on master targets
\texttt{[self-hosted, linux, w33-tpm]}, requires a real TPM2 device and measured
boot log, runs \texttt{tpm2\_makecredential}/\texttt{tpm2\_activatecredential},
requires the enrollment result to state \texttt{hardwareBacked=true}, captures
configured RAPL and/or Kepler evidence into the signed host bundle, and then
requires the portable verifier to pass before uploading an audit artifact.  At
the time of this insert the hardware request had been issued but the self-hosted
job was still queued.  A queued request is not evidence that a physical bundle
exists; only a completed hardware job plus its verifier-approved artifact can
close that claim.

The focused hosted-software closure was likewise queued at publication time.
Accordingly the repository contains the executable certificates and their
fail-closed acceptance logic, while hosted-run numerical metrics remain subject
to the corresponding workflow result rather than being asserted in advance.
